% ============================================================
% Chapter 5: Development
%   Section: The Inference Service
% Included from chapters/05-development/development.tex
% Source: docs/decisions/inference.md.
% ============================================================

\section{The Inference Service}
\label{sec:dev-inference}

Why the three transformer models were moved out of the backend is
explained in Section~\ref{sec:architecture}. This section describes how
the move was made.

\paragraph{Adapters, not a rewrite.} The entity extractor, the
sentiment analyser and the embedding service kept their exact public
signatures and now delegate to an HTTP client. Every call site was
checked before the change: nothing downstream --- the enrichment
pipeline, the claim extractor, the topic classifier, the corrector, the
retrieval and ranking chain --- ever touched a model object, only the
methods that return results. Nothing downstream changed.

\paragraph{A model server that knows nothing of the domain.} The
service loads the models and returns entities, sentiment scores and
vectors. Topics, thresholds and scoring stay in the backend: the topic
classifier still owns the topic vocabulary and the cosine arithmetic,
and only fetches its vectors over HTTP. Tuning therefore stays inside
the backend's configuration and lineage instead of being split across
two deployable units. The three models share one service because they
share a dependency stack, a warm-up and a resource pool.

\paragraph{Failure differs by component, on purpose.} If entity
extraction fails it returns nothing, and the claim scorer, which treats
entities as one weak signal among several, carries on. If sentiment or
embeddings fail they raise, and the article fails: sentiment is a core
input of the admission decision and embeddings feed duplicate
detection, ranking and topic classification, so a silent default would
corrupt decisions without a trace.

\paragraph{Readiness.} The service reports healthy only once all three
models are loaded; the models warm up in a background thread so the
server can answer ``not ready'' honestly instead of blocking its own
start. The compose file holds the backend back until then, which turns
the old first-request stall into one explicit wait at start-up. The
start period is 900~s, measured rather than guessed: a cold model cache
took about eleven minutes to become healthy.

\paragraph{Pinned versions.} The service pins the exact library
versions the backend had resolved. Left open, the resolver chose a new
major version of \texttt{transformers}, which broke GLiNER's tokenizer
loading outright.

\paragraph{Where the CPU time went.} A profile of the service on the
development laptop's CPU (16~October) found three costs that bought
nothing, and one error. The ranker's eight evidence pages were embedded
as one padded batch, which on a CPU took 9.4~s against 6.8~s one at a
time (7.3~s against 4.2~s with the mixed lengths it really sends), with
identical vectors; the service now embeds one text at a time unless
configured for a GPU. Enrichment embedded each article twice, once for
storage and once for the topic classifier, at 3.0~s a time for a
5,700-character article; the classifier now reuses the vector. And the
claim extractor sent one entity request per sentence, 37 for that
article; one batched request finds the same entities in 1.5~s instead
of 2.7~s. The error was GLiNER's input limit: a whole article was
truncated to its first 384 tokens with only a warning, so the
article-level entities came from its opening alone --- 5 found where
its sentences hold 52. Long texts are now read in sentence-bounded
chunks. Together, enrichment of that article went from 17.2~s to
14.1~s while reading all of it.

\paragraph{One trap left by the split.} The backend's configuration
still names the sentiment and embedding models, but the names are
documentary: the models load from the inference service's own
configuration. Only the embedding dimension, which sizes the vector
collection, is still read by the backend, and a test checks it against
the live service.
